% =====================================================================
% LilyTorch: GPU-Accelerated Immersed-Boundary CFD for
% Neuromechanical Simulation of Aquatic Locomotion
%
% IEEE Robotics and Automation Letters (RA-L) submission template.
% Uses IEEEtran with the "journal" option, which is the format
% required by RA-L.
% =====================================================================

\documentclass[journal,twoside,10pt]{IEEEtran}

% ---- Packages --------------------------------------------------------
\usepackage{amsmath,amssymb,amsfonts}
\usepackage{graphicx}
\usepackage{cite}
\usepackage{algorithmic}
\usepackage{array}
\usepackage{booktabs}
\usepackage{url}
\usepackage[hidelinks]{hyperref}
\usepackage{bm}
\usepackage{siunitx}
\usepackage{subcaption}
\usepackage{xcolor}
\usepackage{xspace}

% ---- Convenience macros ---------------------------------------------
\newcommand{\lilytorch}{\textsc{LilyTorch}\xspace}
\newcommand{\farms}{\textsc{FARMS}\xspace}
\newcommand{\mujoco}{\textsc{MuJoCo}\xspace}
\newcommand{\pytorch}{\textsc{PyTorch}\xspace}
\newcommand{\waterlily}{\textsc{WaterLily.jl}\xspace}
\newcommand{\bdim}{BDIM\xspace}

% ---- Paper metadata --------------------------------------------------
\title{\lilytorch: A GPU-Accelerated Immersed-Boundary \\
       Navier--Stokes Solver for Closed-Loop
       Neuromechanical Simulation of Aquatic Locomotion}

\author{Andrea~Ferrario~\IEEEmembership{}%
        \thanks{Manuscript submitted to IEEE Robotics and Automation
        Letters.  The author is with the BioRobotics Laboratory,
        EPFL, Lausanne, Switzerland.
        E-mail: \texttt{andrea.ferrario@epfl.ch}.}%
}

% RA-L running heads
\markboth{IEEE Robotics and Automation Letters. Preprint Version.}%
{Ferrario: \lilytorch --- GPU Immersed-Boundary CFD for Aquatic Locomotion}

\IEEEtitleabstractindextext{%
\begin{abstract}
Simulating freely-swimming robots and animals requires coupling a
multibody dynamics engine with a fluid solver that is fast enough for
closed-loop neuromechanical experiments yet accurate enough to capture
viscous boundary layers and unsteady wakes.  We present \lilytorch, an
open-source, GPU-accelerated incompressible Navier--Stokes solver that
implements a second-order Boundary Data Immersion Method (BDIM2) on a
staggered Cartesian grid and couples two-way with the \mujoco physics
engine through the \farms framework.  The solver supports both
analytical signed-distance functions (spheres, capsules, boxes, fish
spines) and mesh-based bodies (STL/OBJ through \textsc{Open3D} and
fast-marching), handles per-link articulated kinematics, and offers a
choice of spectral (DCT / free-space Green's function) or
variable-coefficient multigrid pressure-projection solvers.  Optional
Smagorinsky large-eddy closure, Carreau / Herschel--Bulkley
non-Newtonian rheology, and a sponge damping layer extend the solver to
a range of bio- and soft-robotic scenarios.  Because the complete
pipeline (advection--diffusion, projection, BDIM, force integration) is
expressed in \pytorch tensors, every hot kernel is JIT-compiled via
\texttt{torch.compile} and runs on a single GPU without domain
decomposition.  We describe the continuous and discrete formulations,
detail how rigid, articulated, and mesh bodies are integrated with
\mujoco, and discuss the implications of the proposed coupling for
embodied studies of aquatic locomotion.
\end{abstract}

\begin{IEEEkeywords}
Biologically-inspired robots, Computational fluid dynamics,
Fluid--structure interaction, Immersed boundary method, Simulation,
Swimming robots.
\end{IEEEkeywords}}

\begin{document}
\maketitle
\IEEEdisplaynontitleabstractindextext
\IEEEpeerreviewmaketitle


% =====================================================================
\section{Introduction}
% =====================================================================
\IEEEPARstart{U}{nderstanding} how neural control, body mechanics, and
the surrounding fluid interact during swimming is a long-standing
question in both neuroscience and robotics.  Anguilliform
swimmers~\cite{tytell2004hydrodynamics}, zebrafish
larvae~\cite{muller2008larval,kern2006simulations}, salamanders
walking and swimming~\cite{ijspeert2007salamander}, and an increasing
number of soft and articulated aquatic
robots~\cite{porez2014improved,katzschmann2018exploration} all exhibit
behaviour that can only be explained by closing the loop between the
controller, the body, and the hydrodynamic
environment~\cite{tytell2010interplay}.  Building computational
platforms that \emph{embody} this loop --- i.e.\ that simulate
controller, multibody dynamics, and fluid together in real- or
near-real-time --- remains a key methodological challenge.

Classical body-conforming Navier--Stokes solvers based on finite volumes
or spectral elements
(e.g.\ \textsc{OpenFOAM}~\cite{weller1998openfoam},
\textsc{Nek5000}~\cite{fischer2007nek5000}) deliver excellent fidelity
but require mesh motion or remeshing when the geometry is articulated
or freely moving; they are also too slow to be driven interactively
from a robotics engine.  Immersed-boundary (IB)
methods~\cite{peskin2002immersed,mittal2005immersed} circumvent these
issues by solving the fluid on a fixed Cartesian grid and representing
the body as a forcing or blending term.  Among IB variants, the
\emph{Boundary Data Immersion Method}
(BDIM)~\cite{weymouth2011boundary,maertens2015accurate} blends the
fluid velocity with the prescribed body velocity through a smoothed
Heaviside function derived from a signed-distance field, and has been
shown to achieve second-order accuracy on a wide range of
bio-inspired flows~\cite{lauder2015fish,bhalla2013unified}.  The
\waterlily package~\cite{weymouth2024waterlily} demonstrated that BDIM
can be implemented concisely on top of modern array frameworks; our
work adopts the same philosophy in \pytorch and extends it to two-way
coupling with a general-purpose multibody engine.

On the robotics side, \mujoco~\cite{todorov2012mujoco} has become the
reference engine for articulated rigid-body dynamics and is widely used
for legged and swimming locomotion.  The
\farms framework~\cite{arreguit2023farms} provides a high-level
neuromechanical abstraction on top of \mujoco, with central-pattern
generators~\cite{ijspeert2008central}, sensor buses, and extension
hooks.  However, neither \mujoco nor \farms natively models a viscous
fluid: aquatic experiments have historically relied on drag
coefficients or added-mass approximations, which cannot reproduce wake
reattachment, vortex shedding, or inter-body interference.  Plugging a
fully-resolved CFD solver into this stack requires:
(i) a fluid formulation that accepts arbitrarily-moving, articulated
geometry at every time step;
(ii) a body representation that covers analytical primitives as well
as arbitrary triangular meshes (for which \mujoco uses STL/OBJ); and
(iii) a pressure-projection path that tolerates per-step changes in
body mask, density, and viscosity.

This paper describes \lilytorch, an open-source
\pytorch-based~\cite{paszke2019pytorch} incompressible Navier--Stokes
solver that we designed to fill this gap.  Its contributions are:

\begin{itemize}
\item \textbf{A single-GPU, two-way coupled CFD+\mujoco pipeline.}
  Every time step the fluid solver reads link poses from \mujoco,
  recomputes per-link signed-distance functions (SDFs), performs a
  BDIM2 Heun predictor--corrector step, integrates hydrodynamic
  forces, and writes them back as external wrenches.  The entire
  pipeline runs on one GPU without inter-process communication.
\item \textbf{Unified analytical / mesh body handling.}  A common
  \texttt{Body} abstraction exposes spheres, capsules, boxes, NACA
  fish-spine profiles, and STL/OBJ meshes (loaded via
  \textsc{Open3D}~\cite{zhou2018open3d} with fast-marching sign
  propagation~\cite{sethian1996fastmarch}) through the same BDIM
  interface.  Articulated animats are supported through composite
  bodies that take the pointwise union of per-link SDFs.
\item \textbf{Choice of pressure solvers.}
  A spectral DCT solver (Neumann boundaries), a free-space Green's
  function convolution~\cite{hejlesen2013high} for unbounded domains,
  and a geometric multigrid~\cite{trottenberg2001multigrid}
  preconditioned CG (MGCG) solver for variable-coefficient problems
  arising from the BDIM mask or variable density.
\item \textbf{Advection, turbulence, and rheology models.}  Six
  advection schemes (QUICK~\cite{leonard1979quick},
  ADBQUICKEST, CUBISTA~\cite{alves2003convergent},
  van~Leer~\cite{vanleer1979towards}, CDS, and semi-Lagrangian
  backtracing~\cite{stam1999stable}); an optional
  Smagorinsky~\cite{smagorinsky1963general} large-eddy closure;
  Carreau and Herschel--Bulkley non-Newtonian
  models~\cite{bird1987dynamics}; and a quadratic sponge layer for
  open-domain emulation.
\item \textbf{GPU kernel fusion.}  All hot paths
  (advection--diffusion, multigrid smoother and V-cycle, force
  integration, SDF evaluation) are expressed as pure functions that
  can be JIT-fused through \texttt{torch.compile}, taking advantage
  of CUDA-graph capture.
\end{itemize}

The remainder of this paper is organised as follows.
Section~\ref{sec:methods} presents the continuous equations, the
discrete numerical schemes, the BDIM2 immersed-boundary formulation,
the analytical and mesh body classes, and the coupling with \mujoco
through \farms.  Section~\ref{sec:discussion} discusses design
trade-offs, comparisons with existing tools, limitations, and
planned extensions.  Experimental validation and quantitative results
are deferred to a companion section.


% =====================================================================
\section{Methods}\label{sec:methods}
% =====================================================================

\subsection{Governing equations}\label{sec:governing}

\lilytorch solves the incompressible Navier--Stokes equations in
primitive variables,
\begin{align}
  \frac{\partial \mathbf u}{\partial t}
  + (\mathbf u\!\cdot\!\nabla)\mathbf u
  &= -\frac{1}{\rho}\nabla p + \nu\,\nabla^{2}\mathbf u,
  \label{eq:momentum}\\[2pt]
  \nabla\!\cdot\!\mathbf u &= 0,
  \label{eq:continuity}
\end{align}
where $\mathbf u=(u,v)$ in 2-D or $(u,v,w)$ in 3-D, $p$ is the
pressure, $\rho$ the (constant) fluid density, and $\nu$ the kinematic
viscosity.  The incompressibility constraint~\eqref{eq:continuity} is
enforced by the classical Chorin--T\'emam fractional-step
method~\cite{chorin1968numerical,temam1969approximation}:
an explicit provisional velocity $\tilde{\mathbf u}$ is first advanced
with the momentum terms, and a Poisson problem is then solved to
obtain the pressure correction.


\subsection{Staggered MAC grid and time integration}\label{sec:grid}

The domain $\Omega\!\subset\!\mathbb R^{d}$ ($d\!\in\!\{2,3\}$) is
discretised on a Marker-and-Cell (MAC) staggered
grid~\cite{harlow1965numerical}: pressure is stored at cell centres
and each velocity component on its own face grid
($u$ at $x$-faces, $v$ at $y$-faces, $w$ at $z$-faces).  One ghost
cell is added on each side of $\Omega$, giving array shapes
$(N_x{+}2)\!\times\!(N_y{+}2)\,[\times(N_z{+}2)]$; boundary conditions
are enforced through image/mirror ghost values for Dirichlet or
zero-gradient ghost values for Neumann faces.

Time integration follows Heun's method (explicit RK2) as in
\waterlily~\cite{weymouth2024waterlily}:
\begin{enumerate}
\item \textbf{Predictor:} advection--diffusion of $\mathbf u^{n}$
  gives $\mathbf u^{*}$; BDIM blends $\mathbf u^{*}$ with the
  body velocity (Section~\ref{sec:bdim}); pressure projection with
  weight $w{=}1$ produces $\tilde{\mathbf u}$.
\item \textbf{Corrector:} advection--diffusion of $\tilde{\mathbf u}$
  gives $\mathbf u^{**}$; a half-step rebase
  $\mathbf u^{**}\!\leftarrow\!\tfrac12(\mathbf u^{**}+\mathbf u^{n})$
  is applied; BDIM is re-blended; a second pressure projection with
  weight $w{=}\tfrac12$ returns $\mathbf u^{n+1}$.
\end{enumerate}
The advection--diffusion sub-step itself uses forward Euler.  CFL
stability requires
$\Delta t \le \min(h)\,/\,(v_{\max} + 3(\nu+\nu_t^{\max}))$,
where $\nu_t$ is the (optional) Smagorinsky eddy viscosity.

\paragraph{Advection schemes}  Six choices are exposed for the flux
$\nabla\!\cdot\!(\mathbf u\!\otimes\!\phi)$:
third-order upwind QUICK~\cite{leonard1979quick} with median limiter;
a TVD-limited variant (ADBQUICKEST); the high-resolution
CUBISTA~\cite{alves2003convergent}; van~Leer's
limiter~\cite{vanleer1979towards}; plain central differences; and
Stam's unconditionally stable semi-Lagrangian
back-tracing~\cite{stam1999stable} for very large $\Delta t$ at the
cost of numerical diffusion.

\paragraph{Optional Smagorinsky LES}  When the grid is too coarse to
resolve all eddies, the subgrid-scale closure of
Smagorinsky~\cite{smagorinsky1963general} adds an eddy viscosity
$\nu_t=(C_s\Delta)^{2}|\bar S|$, with
$|\bar S|=\sqrt{2\,\bar S_{ij}\bar S_{ij}}$ the resolved strain-rate
magnitude and $\Delta=h$.  The momentum diffusion becomes the
variable-coefficient operator
$\nabla\!\cdot\!\bigl[(\nu+\nu_t)\nabla\mathbf u\bigr]$,
discretised by arithmetic averaging of $(\nu+\nu_t)$ at the MAC
faces.  Setting $C_s{=}0$ recovers the constant-viscosity path at
zero overhead.


\subsection{Boundary Data Immersion Method}\label{sec:bdim}

Bodies are represented by a signed-distance function (SDF)
$d(\mathbf x)$, with $d>0$ in the fluid, $d=0$ on the surface, and
$d<0$ inside the body.  Two smooth kernel functions of half-width
$\varepsilon{=}1.5h$ are built from $d$: a Heaviside-like indicator
\begin{equation}
  \mu_0(d)\!=\!
  \begin{cases}
    0,& d\!\le\!-\varepsilon,\\
    \tfrac12\!\left[1+\tfrac{d}{\varepsilon}
      +\tfrac{\sin(\pi d/\varepsilon)}{\pi}\right],
      & |d|\!<\!\varepsilon,\\
    1,& d\!\ge\!\varepsilon,
  \end{cases}
  \label{eq:mu0}
\end{equation}
and a first-moment kernel $\mu_1(d)$ obtained by integrating
$\mu_0$.  At each RK2 sub-step the provisional velocity $\phi$ is
blended with the prescribed body velocity $\mathbf v_b$ via the
BDIM2 \emph{meta-equation}~\cite{maertens2015accurate,weymouth2011boundary}:
\begin{equation}
  \phi_\mathrm{out}
  = \mu_0\,(\phi-\mathbf v_b)
  + \mathbf v_b
  + \mu_1\,\hat{\mathbf n}\!\cdot\!\nabla(\phi-\mathbf v_b),
  \label{eq:bdim}
\end{equation}
with $\hat{\mathbf n}\!=\!\nabla d/|\nabla d|$.  Inside the body
$\mu_0,\mu_1\!\to\!0$ and $\phi_\mathrm{out}\!=\!\mathbf v_b$; in the
fluid $\mu_0\!\to\!1,\mu_1\!\to\!0$ and $\phi$ is recovered; the
transition layer provides second-order enforcement of the no-slip
boundary condition without sub-grid reconstruction.

The pressure-projection step is likewise modified to respect the
body mask,
\begin{equation}
  \nabla\!\cdot\!\!\left(\frac{w\,\Delta t}{\rho}\,\mu_0\,\nabla p\right)
  = \nabla\!\cdot\!\tilde{\mathbf u},
  \label{eq:poisson}
\end{equation}
so that the pressure gradient vanishes inside the body and the
projection remains well-posed at the surface.

\paragraph{Force integration}
Hydrodynamic forces on a body are recovered by integrating the
Cauchy stress tensor over a smoothed surface using an $\varepsilon$-
width mollified Dirac:
\begin{align}
  \mathbf F_\mathrm{visc}
  &=\int(\bm\sigma\!\cdot\!\hat{\mathbf n})\,
      \delta_\varepsilon(d-\varepsilon)\,\mathrm dV,\\
  \mathbf F_\mathrm{pres}
  &=-\int p\,\hat{\mathbf n}\,
      \delta_\varepsilon(d)\,\mathrm dV,
\end{align}
with
$\delta_\varepsilon(d)=[1+\cos(\pi d/\varepsilon)]/(2\varepsilon)$.
Torques are accumulated directly from component sums to avoid
allocating full moment-arm tensors on the GPU.


\subsection{Pressure-projection solvers}\label{sec:poisson}

\lilytorch provides three Poisson back-ends.

\paragraph{FFT / DCT (Neumann)}  When every face carries
$\partial p/\partial n{=}0$, the discrete Neumann Laplacian is
diagonalised by the type-II Discrete Cosine
Transform~\cite{makhoul1980dct},
$\lambda_k=(2/h^2)[\cos(\pi k/N)-1]$, and the solve reduces to
DCT--divide--iDCT.  This is the fastest path and is used for the
majority of validation benchmarks.

\paragraph{Free-space Green's function}  For unbounded domains the
RHS is zero-padded to $2N$ and convolved with a regularised Green's
function.  In 2-D we use the eighth-order algebraic smoothing of
Hejlesen et al.~\cite{hejlesen2013high}; in 3-D the Gaussian/erf
regularisation $G(r)=-(4\pi r)^{-1}\mathrm{erf}(r/\sqrt 2\sigma)$
with $\sigma{=}h$.  Green's functions are precomputed once and cached
on disk.

\paragraph{Geometric multigrid / MGCG}  When the BDIM mask or a
variable fluid/body density makes~\eqref{eq:poisson} genuinely
variable-coefficient, we solve it with a V-cycle
multigrid~\cite{trottenberg2001multigrid} wrapped, if requested,
inside a conjugate-gradient iteration (MGCG).  The discrete stencil
along dimension $d$ reads
\begin{equation*}
  \frac{c_{d+}p_{i+1}-(c_{d+}+c_{d-})p_{i}+c_{d-}p_{i-1}}{h^{2}},
\end{equation*}
with face coefficients $c_{d\pm}$ taken as the harmonic average of
$\mu_0$ (or $\Delta t/\rho$ in the variable-density case) at cell
faces.  Available smoothers are weighted Jacobi
($\omega{=}0.7$ default) and red/black Gauss--Seidel; full-weighting
restriction and bilinear (2-D) or trilinear (3-D) prolongation are
used.  The coarsest level is a trivial $2\times 2\,(\times 2)$ direct
solve.

\paragraph{Variable-density extension for \mujoco coupling}  When the
body density $\rho_b$ differs from the fluid $\rho_f$ we assemble an
effective density
$\rho(\mathbf x)=\rho_b+(\rho_f-\rho_b)\mu_0(\mathbf x)$ and face
coefficients $c=\Delta t/\rho$ directly on the MAC faces.  This
coupled variable-density operator departs from the single-density
BDIM form of \waterlily and is therefore solved exclusively by the
multigrid path.


\subsection{Body representation: analytical vs mesh}\label{sec:bodies}

The abstract \texttt{Body} class exposes three operations:
(i) \texttt{compute\_sdf} evaluates $d$ on the (possibly staggered)
grid; (ii) \texttt{compute\_velocity} returns $\mathbf v_b$; and
(iii) \texttt{update(t, it)} advances the body kinematics.  Two
concrete families are provided.

\subsubsection{Analytical bodies}  For primitives whose SDF admits a
closed form we evaluate $d$ directly on the staggered grid at each
step.  Built-in shapes are spheres / circles, axis-aligned boxes,
capsules (useful for articulated fish tails), and NACA-foil fish
spines parameterised by a cubic spline of centre-line
positions~\cite{kern2006simulations}.  Because $d$ is evaluated
analytically at every time step, arbitrary kinematics (prescribed or
driven by \mujoco) cost only a small number of elementwise GPU
kernels.

\subsubsection{Mesh bodies}  For arbitrary
triangulated geometry (STL/OBJ) the SDF is computed in a
pre-processing pass and then transformed rigidly at run time.  The
pipeline is: closest-point queries through
\textsc{Open3D}~\cite{zhou2018open3d} ray-casting; sign determination
via winding numbers; narrow-band fast
marching~\cite{sethian1996fastmarch} (scikit-fmm) to extend the
distance field beyond the mesh; and projection onto the three MAC
face grids.  At simulation time the body pose (from \mujoco or from
a prescribed trajectory) rigidly transforms the cached SDF via a
\texttt{pytorch-interpolation} regular-grid trilinear interpolator,
so that mesh updates are a single GPU interpolation rather than a
re-triangulation.

\subsubsection{Composite and articulated bodies}
The union of two bodies is represented through the pointwise minimum
of their SDFs, $d_\mathrm{union}=\min(d_1,d_2)$.  A
\texttt{CompositeSegmentBody} stacks the per-link SDFs of an
articulated animat, while \texttt{MultiAnimatBodies} aggregates
independent swimmers in the same domain.  An axis-aligned
bounding-box optimisation avoids evaluating each link's SDF over the
full grid whenever it occupies less than $\approx 90\%$ of the
domain, which is the common case for small animats in large pools.


\subsection{Coupling with \mujoco through \farms}\label{sec:coupling}

In standalone mode the \texttt{FluidSolver} simply advances
Eqs.~\eqref{eq:momentum}--\eqref{eq:continuity} together with
prescribed body kinematics.  In coupled mode, the solver is invoked
as a \farms \texttt{TaskExtension} (\texttt{FluidExtension}), which
participates in \mujoco's \texttt{before\_step} callback.  One
coupled time step proceeds as follows:
\begin{enumerate}
\item \emph{Pose read-out.}  \farms sensors expose link positions,
      orientations, and linear/angular velocities; these are copied
      into \pytorch tensors.
\item \emph{SDF assembly.}  Each link updates its cached SDF
      (analytical re-evaluation for primitives, rigid-transform
      interpolation for meshes) inside a narrow AABB; the union
      over links is taken via $\min$.
\item \emph{Fluid step.}  The Heun predictor--corrector advances the
      velocity field, applying BDIM at every RK2 stage and solving
      the pressure Poisson equation on the current mask.
\item \emph{Force integration.}  Viscous and pressure contributions
      are integrated per link, plus a hydrostatic buoyancy term
      $F_{z,\mathrm{buoy}}=-\rho_w (m_\ell/\rho_b) g\,
      \min\!\bigl((z_\mathrm{surf}+h_\ell-z_\ell)/(2h_\ell),1\bigr)$
      that smoothly switches on as a link crosses a prescribed free
      surface.
\item \emph{Wrench write-back.}  The six-vector per-link wrench is
      written to \mujoco's \texttt{xfrc\_applied} buffer and
      integrated by the rigid-body engine during its next step.
\end{enumerate}
Because \lilytorch and \mujoco share a single Python process, no
inter-process communication is required; the GPU and CPU execute in
alternation once per simulation step.  A complementary
\texttt{WaterDynamicsCallback} feeds the CFD velocity field back to
\mujoco's native drag solver for cheaper, lower-fidelity scenarios
(e.g.\ shallow-water robots), and a \texttt{FlowViewer} extension
renders vorticity isosurfaces directly inside the \mujoco GUI and
recorded video, aiding debugging.

\paragraph{Body density and buoyancy}  Hydrodynamic and inertial
properties of the multibody model are inherited from the \mujoco XML
(masses, inertias, collision meshes).  The fluid solver additionally
queries a \emph{body density} $\rho_b$: when it equals the fluid
density the projection reduces to the standard BDIM form; otherwise
the variable-density multigrid path is taken and buoyancy is computed
from the displaced volume.  For mesh-based animats (e.g.\ the
pleurodeles and salamander models shipped with the repository) the
displaced volume is tabulated once from the SDF and reused; for
analytical composites it is the sum of the primitive volumes.


\subsection{Non-Newtonian rheology and sponge layer}\label{sec:extras}

Two additional physical models extend the solver to soft-robotic and
open-domain scenarios.  A Carreau fluid
model~\cite{bird1987dynamics}
$\nu(\dot\gamma)=\nu_\infty+(\nu_0-\nu_\infty)[1+(\lambda\dot\gamma)^2]^{(n-1)/2}$
captures shear-thinning behaviour (e.g.\ swimming in
carboxymethyl-cellulose gels), and a Herschel--Bulkley extension
adds a yield-stress term
$\tau_y/[\rho\max(\dot\gamma,\epsilon)]$ that is clamped from above
by the CFL diffusion limit.  A quadratic sponge layer of thickness
$L_s$ and strength $\sigma_\mathrm{max}$ damps outgoing disturbances
through the post-projection update
$\mathbf u\!\leftarrow\!\mathbf u/(1+\Delta t\,\sigma)$, which
emulates an infinite far-field at a fraction of the cost of
perfectly-matched layers.


\subsection{Implementation and GPU kernel fusion}\label{sec:impl}

The solver is written in approximately 10\,k lines of Python on top
of \pytorch~\cite{paszke2019pytorch}.  All tensors live on a single
CUDA device; multi-GPU domain decomposition is not required for the
problem sizes considered here ($N_x N_y N_z\!\lesssim\!256^{3}$
typical, $512^3$ achievable).  Hot paths --- the advection--diffusion
kernel, the multigrid smoother and full V-cycle, the SDF and
$\mu_0$/$\mu_1$ evaluation, and the force integration --- are
expressed as \emph{module-level pure functions} rather than methods,
so that \texttt{torch.compile} with \texttt{mode="reduce-overhead"}
can fuse them and capture CUDA graphs.  Variants are generated per
spatial dimension and per smoother choice.

YAML configuration files drive every simulation: the
\texttt{solver}, \texttt{boundary\_conditions}, \texttt{body}, and
\texttt{output} blocks select grid, BCs, bodies, and logging; an
additional \texttt{physics} block is used only when coupling with
\farms.  State checkpoints are written to HDF5 and can be resumed.


% =====================================================================
\section{Discussion}\label{sec:discussion}
% =====================================================================

\subsection{Design trade-offs}

\lilytorch occupies a middle ground in the landscape of
immersed-boundary solvers.  Highly-optimised C++/CUDA codes such as
\textsc{IBAMR}~\cite{bhalla2013unified} and
\textsc{AMReX}-based solvers offer adaptive mesh refinement and
distributed-memory scaling, at the cost of a steep learning curve
and a much less flexible body
interface.  \waterlily~\cite{weymouth2024waterlily} pioneered the
idea of a BDIM solver that is short enough to be re-read and
hackable; our work preserves that philosophy but moves the
implementation to \pytorch, which grants three practical
advantages.  First, \pytorch's autograd and \texttt{torch.compile}
infrastructure are readily available for \emph{adjoint}-based shape
and control optimisation~\cite{kochkov2021machinelearned}, without
any additional code.  Second, the Python ecosystem makes it trivial
to embed the solver inside \mujoco, \farms, or reinforcement-learning
training loops, which is our primary use case.  Third, every
research group already familiar with \pytorch can modify the solver
--- for instance by swapping the advection scheme or plugging in a
learned closure~\cite{brunton2020machine} --- without cross-language
FFI.

The price paid is a single-GPU ceiling: with $512^3$ grids at single
precision we are memory-bound on a 40\,GB accelerator and cannot
(yet) scale out.  For the 3-D experiments reported in the companion
section (free-swimming anguilliform animats, salamander paddling,
single-sphere sedimentation benchmarks against
Koumoutsakos \& Leonard~\cite{koumoutsakos1995high} and
Bergmann \& Iollo~\cite{bergmann2011accurate}) this ceiling is
comfortable; beyond it, either domain decomposition or the existing
PETSc back-end~\cite{balay2019petsc} is the most promising path.


\subsection{Comparison with analytical drag models}

Most robotics-oriented swimming simulators approximate hydrodynamics
through link-wise drag and added-mass
tensors~\cite{porez2014improved}, which are computationally cheap but
discard wake-induced effects, inter-link shielding, and boundary
corrections.  By solving the full Navier--Stokes equations on a
staggered MAC grid with BDIM2, \lilytorch reproduces phenomena that
are structurally absent from drag-based pipelines: vortex shedding
behind a pinned cylinder (validated against
Koumoutsakos \& Leonard~\cite{koumoutsakos1995high}), thrust
enhancement from body undulation, and wake--wake interaction between
swimmers.  The cost is roughly two to three orders of magnitude
higher wall-clock time per step than a pure drag simulator, which is
acceptable for scientific studies (hours to days per trajectory) but
precludes, for the moment, training full reinforcement-learning
policies at scale unless low-resolution proxies are used.


\subsection{Body-representation considerations}

The dual analytical/mesh body path is crucial for embodied robotics
experiments.  Analytical primitives keep the SDF computation
tractable when an articulated animat has dozens of links whose
shapes would otherwise be expensive to interpolate every step; they
also sidestep the mesh-resolution / grid-resolution mismatch that
can introduce high-frequency noise at the BDIM boundary.  Conversely,
mesh bodies are indispensable whenever the geometry comes from
imaging data (fish micro-CT, salamander body scans) or has been
optimised offline.  Our use of \textsc{Open3D} ray-casting, winding
numbers, and fast marching yields an SDF of quality comparable to
signed-distance generators built on top of
\textsc{CGAL}~\cite{alliez2009cgal}, but in pure Python and without
C++ toolchain dependencies --- a practical advantage for a library
that must be installable on a wide variety of robotics workstations.
The rigid-transform interpolation performed at run time preserves
single-GPU throughput: a typical STL body costs only a few extra
milliseconds per step, independent of the number of triangles.


\subsection{Coupling stability and time-step considerations}

Two-way coupling with \mujoco exposes two numerical pitfalls.
\emph{Added-mass instability} can arise when $\rho_b\!\approx\!\rho_f$
and the fluid forces dominate the body inertia~\cite{causin2005added};
for the neutrally-buoyant animats studied here we mitigate it with
the variable-density projection
$\rho=\rho_b+(\rho_f-\rho_b)\mu_0$, which absorbs the fluid's
virtual mass into the Poisson operator rather than treating it as an
explicit body force.  \emph{Temporal resolution mismatch}
between the fluid (CFL-limited at $\Delta t\!\sim\!10^{-4}$~s on a
$1$~cm grid) and \mujoco (typically $10^{-3}$~s for contact
stability) is handled by running both engines at the fluid time
step; the overhead on the rigid-body side is negligible compared to
the fluid cost.

Finally, the sponge layer (Section~\ref{sec:extras}) is particularly
important for low-Reynolds swimmers whose pressure Poisson equation
instantaneously communicates the body motion to the whole domain; a
few-centimetre layer at the walls suppresses spurious recirculation
without contaminating the near-field, as verified on single-sphere
sedimentation benchmarks.


\subsection{Limitations and outlook}

At present \lilytorch is limited to uniform Cartesian grids and
single-phase flow; free-surface extensions through a volume-of-fluid
or level-set approach, and local grid refinement through an AMReX
back-end, are natural next steps.  Adaptive time stepping is
currently manual and could be automated from the flow diagnostics
already tracked by the solver (kinetic energy, enstrophy, max-CFL,
max-divergence).  On the robotics side, a tighter integration of the
fluid forces into \mujoco's semi-implicit integrator would address
the residual added-mass instability for fully hollow bodies.  Finally,
because all kernels are \pytorch-native and differentiable, end-to-end
adjoint optimisation of gait parameters or morphology appears
feasible~\cite{kochkov2021machinelearned,brunton2020machine} and will
be the subject of follow-up work.


% =====================================================================
% Acknowledgments (placeholder)
% =====================================================================
\section*{Acknowledgments}
The author thanks the members of the BioRobotics Laboratory for
discussions and for providing the experimental datasets used in the
companion validation section.


% =====================================================================
\bibliographystyle{IEEEtran}
\bibliography{references}
% =====================================================================

\end{document}
